%% if A
\chapter{Alle CoolProp-opzoekingen op een rij}
\label{b:opzoek}
Dit is de lijst van wat CoolProp voor de visuele PFD berekent. Per opzoeking staat welke twee eigenschappen de toestand
vastleggen (toestandspostulaat) en wat je in Cengel zou doen. Eenheden zoals in het verslag; CoolProp zelf rekent in K, Pa,
J/kg.
\begin{center}\scriptsize\setlength{\tabcolsep}{3pt}
\begin{longtable}{p{3.2cm}p{1.6cm}p{3.4cm}rp{3.4cm}}
\toprule
wat & stof & invoer (2 eigenschappen) & uitkomst & Cengel-equivalent\\
\midrule\endhead
dauwpunt S13 & Water & $P=\VAR{n(R.P_v,3)}$\,kPa, $x=1$ & \VAR{n(R.T_dauw,2)}\textdegree C & A-5 of interpoleren in A-4\\
$P_\mathrm{sat}$ KO-drum & Water & $T=30$\textdegree C, $x=0$ & \VAR{n(R.Psat_KO,4)}\,kPa & A-4: 4,2469\\
$h$ \OO\ 70\textdegree C & Oxygen & $T$, $p_{\OO}=\VAR{n(R.pO2_in,1)}$\,kPa & \VAR{n(R.hO2_70,2)}\,kJ/kg & ideaal gas, $c_p=0{,}918$\\
$h$ \OO\ 30\textdegree C & Oxygen & $T$, $p_{\OO}=\VAR{n(R.pO2_ko,1)}$\,kPa & \VAR{n(R.hO2_30,2)}\,kJ/kg & idem\\
$h_v$ 70 / 30\textdegree C & Water & $T$, $P=0{,}5$\,kPa & \VAR{n(R.hv_70,1)} / \VAR{n(R.hv_30,1)} & A-4: $h_g$ 2626,1 / 2555,6\\
$h_f$ 30\textdegree C & Water & $T$, $x=0$ & \VAR{n(R.hf_30,2)}\,kJ/kg & A-4: 125,74\\
$s$, $h$ mengsel K-102 & Oxygen + Water & $T$, partiële drukken & $s_1=\VAR{n(R.s1,4)}$ & ideaal gas, $k=1{,}395$\\
$T_{2s}$, $T_2$ K-102 & idem & $s=s_1$ / $h=h_2$, $P_2$ (bisectie) & \VAR{n(R.T2s,1)} / \VAR{n(R.T_uit_K102,1)}\textdegree C & $T_1 r_p^{(k-1)/k}$\\
dauwpunt S18 & Water & $P=\VAR{n(R.pv_2,2)}$\,kPa, $x=1$ & \VAR{n(R.T_dauw_S18,1)}\textdegree C & A-4 interpoleren\\
$h_1$, $s_1$ K-101 & Hydrogen & $T=25$\textdegree C, $P=\VAR{n(R.P_in_K101,0)}$\,kPa & \VAR{n(R.h1h,1)}; \VAR{n(R.s1h,3)} & ideaal gas\\
$h_{2s}$, $T_{2s}$ K-101 & Hydrogen & $P=\VAR{n(R.P_uit_K101,0)}$\,kPa, $s=s_1$ & \VAR{n(R.h2sh,1)}; \VAR{n(R.T2sh,1)}\textdegree C & $T_1 r_p^{(k-1)/k}$\\
$T_2$ K-101 & Hydrogen & $P_2$, $h=h_2$ & \VAR{n(R.T_uit_K101,2)}\textdegree C & $T_1 + (T_{2s}-T_1)/\eta$\\
$h$ \HH\ bij S6/50/25\textdegree C & Hydrogen & $T$, $P_2$ & \VAR{n(R.hH_S6,1)} / \VAR{n(R.hH_50,1)} / \VAR{n(R.hH_25,1)} & $c_p\Delta T$, $c_p = 14{,}307$\\
$\rho$ V-101 max/min & Hydrogen & $T=25$\textdegree C, $P$ & \VAR{n(R.rho_V101_max,3)} / \VAR{n(R.rho_V101_min,3)} & $P/(ZRT)$\\
$\rho$ V-102 max/min & Hydrogen & $T=25$\textdegree C, $P$ & \VAR{n(R.rho_V102_max,3)} / \VAR{n(R.rho_V102_min,3)} & $P/(ZRT)$\\
$k = c_p/c_v$ & Hydrogen & $T=25$\textdegree C, 1\,atm & \VAR{n(R.k_H2,4)} & A-2: 1,405\\
$v_f$ 30\textdegree C & Water & $T$, $x=0$ & \VAR{n(R.v_f,6)} & A-4: 0,001004\\
$\Delta h$ koelwater 20$\rightarrow$30 & Water & $T$, $P=3$\,bar & \VAR{n(R.dh_kw,3)} & $4{,}18\cdot10$\\
$h$ netwater 40/45/70 & Water & $T$, $P=10$\,bar & \VAR{n(R.hw[40],2)} / \VAR{n(R.hw[45],2)} / \VAR{n(R.hw[70],2)} & $4{,}18\,T$\\
WP-201 punt 1 & \VAR{wp1.fluid} & $T=40$\textdegree C, $x=1$ & $h=\VAR{n(wp1.h1,2)}$, $P=\VAR{n(wp1.Pe,1)}$ & A-11 (R-134a) / diagram\\
WP-201 punt 2s & \VAR{wp1.fluid} & $P_c$, $s=s_1$ & $h=\VAR{n(wp1.h2s,2)}$ & A-13 / diagram\\
WP-201 punt 2 & \VAR{wp1.fluid} & $P_c$, $h=h_2$ & $T=\VAR{n(wp1.T2,2)}$\textdegree C & A-13 / diagram\\
WP-201 punt 3 & \VAR{wp1.fluid} & $T=75$\textdegree C, $x=0$ & $h=\VAR{n(wp1.h3,2)}$, $P=\VAR{n(wp1.Pc,1)}$ & A-11 / diagram\\
WP-201 punt 4 & \VAR{wp1.fluid} & $P_e$, $h=h_3$ & $x=\VAR{n(wp1.x4,4)}$ & $x=(h-h_f)/h_{fg}$\\
WP-202 punt 1 & \VAR{wp2.fluid} & $T=40$\textdegree C, $x=1$ & $h=\VAR{n(wp2.h1,2)}$, $P=\VAR{n(wp2.Pe,1)}$ & diagram\\
WP-202 punt 2s & \VAR{wp2.fluid} & $P_c$, $s=s_1$ & $h=\VAR{n(wp2.h2s,2)}$ & diagram\\
WP-202 punt 2 & \VAR{wp2.fluid} & $P_c$, $h=h_2$ & $T=\VAR{n(wp2.T2,2)}$\textdegree C & diagram\\
WP-202 punt 3 & \VAR{wp2.fluid} & $T=100$\textdegree C, $x=0$ & $h=\VAR{n(wp2.h3,2)}$, $P=\VAR{n(wp2.Pc,1)}$ & diagram\\
WP-202 punt 4 & \VAR{wp2.fluid} & $P_e$, $h=h_3$ & $x=\VAR{n(wp2.x4,4)}$ & $x=(h-h_f)/h_{fg}$\\
$\rho$ koudemiddel R1–R8 & beide & ($T$,$x$) of ($P$,$h$) & zie hoofdstuk~\ref{h:leiding} & isochoren in diagram\\
$\rho$ water/\HH\ leidingen & Water / Hydrogen & $T$, $P$ & zie hoofdstuk~\ref{h:leiding} & $1/v_f$ / $P/(RT)$\\
molaire massa's & O\textsubscript{2}, H\textsubscript{2}O, H\textsubscript{2} & -- & \VAR{n(R.M_O2,4)}; \VAR{n(R.M_H2O,4)}; \VAR{n(R.M_H2,4)} & A-1: 31,999; 18,015; 2,016\\
\bottomrule
\end{longtable}
\end{center}
%% else
\chapter{Alle Cengel-tabelwaarden op een rij}
\label{b:opzoek}
\begin{center}\small
\begin{longtable}{p{6cm}rp{2cm}l}
\toprule
grootheid & waarde & eenheid & tabel\\
\midrule\endhead
$M_{\OO}$, $M_{\HHO}$, $M_{\HH}$ & 31,999; 18,015; 2,016 & kg/kmol & A-1\\
$R_u$ & 8,314 & kJ/(kmol$\cdot$K) & A-1\\
$R_{\HH}$ & 4,124 & kJ/(kg$\cdot$K) & A-1\\
$c_p$, $k$ \OO\ (300\,K) & 0,918; 1,395 & & A-2\\
$c_p$, $k$ \HH\ (300\,K) & 14,307; 1,405 & & A-2\\
$c_p$ waterdamp (300\,K) & 1,8723 & kJ/(kg$\cdot$K) & A-2\\
$c_p$ vloeibaar water & 4,18 & kJ/(kg$\cdot$K) & A-3\\
$P_\mathrm{sat}$ 30 / 65 / 70\textdegree C & 4,2469 / 25,043 / 31,202 & kPa & A-4\\
$h_g$ 30 / 70\textdegree C & 2555,6 / 2626,1 & kJ/kg & A-4\\
$h_f$ 30\textdegree C & 125,74 & kJ/kg & A-4\\
$v_f$ 30\textdegree C & 0,001004 & m$^3$/kg & A-4\\
HHV waterstof & 141\,800 & kJ/kg & A-27\\
tweede viriaalcoëfficiënt \HH\ (25\textdegree C) & 14,8 & cm$^3$/mol & literatuurwaarde\\
\midrule
\multicolumn{4}{l}{\emph{afgelezen uit het log $p$–$h$-diagram (figuren hoofdstuk~\ref{h:wp})}}\\
%% for c, t in ((wp1, "WP-201"), (wp2, "WP-202"))
\VAR{t} (\VAR{c.fluid}): $h_1$; $h_{2s}$; $h_3$ & \VAR{n(c.h1,1)}; \VAR{n(c.h2s,1)}; \VAR{n(c.h3,1)} & kJ/kg & diagram\\
\VAR{t}: $p_e$; $p_c$ & \VAR{n(c.Pe,0)}; \VAR{n(c.Pc,0)} & kPa & diagram\\
\VAR{t}: $h_f(T_e)$; $h_g(T_c)$ & \VAR{n(c.hf_e,1)}; \VAR{n(c.hg_c,1)} & kJ/kg & diagram\\
\VAR{t}: $\rho_1$; $\rho_2$; $\rho_3$; $\rho_4$ & \VAR{g(c.rho1,2)}; \VAR{g(c.rho2,2)}; \VAR{g(c.rho3,3)}; \VAR{g(c.rho4,2)} & kg/m$^3$ & isochoren\\
%% endfor
\bottomrule
\end{longtable}
\end{center}
%% endif
